% ==========================================================================
% GENERATED by scripts/split_paper.py. Do not edit this file.
%
% Source: imported_papers/grokking/iclr2027_conference_v2.tex
% Part 5: The rules that grok share the algorithm but not the circuit
%
% Edits belong upstream, in the published paper. To change what the
% chapter includes, change imports/grokking.toml and re-run
% the splitter.
% ==========================================================================
\section{The rules that grok share the algorithm but not the circuit}
\label{sec:select}

The rules that grok all reach the same test accuracy, and they all implement the same algorithm. What differs is the circuit, the way that algorithm is laid out across the units.

\subsection{Same algorithm, different circuits}
Comparing circuits needs a measure of how many units carry the computation. Indexing the $n$ units of a hidden layer by $j$ and writing $a_j = \|\mW^{\mathrm{out}}_{:,j}\|_2$ for unit $j$'s \emph{outgoing norm}, the norm of the weights carrying its output onward, the \emph{participation ratio} is
\begin{equation}
\mathrm{PR} \;=\; \frac{\big(\sum_j a_j\big)^2}{n \sum_j a_j^2} \;\in\; (0, 1],
\label{eq:pr}
\end{equation}
so $\mathrm{PR} = 1$ when every unit contributes equally, $\mathrm{PR} = 1/n$ when one unit does all the work, and $n \cdot \mathrm{PR}$ reads as an \emph{effective unit count} (Proposition~\ref{prop:pr}). Computed on the outgoing norms we call it \emph{circuit participation}. It does not depend on a choice of basis and does not move when the readout is rescaled (Appendix~\ref{app:pr}).

Backpropagation and DFC sit at opposite extremes. Backpropagation carries the computation on roughly $154$ of $256$ effective units ($n \cdot \mathrm{PR}$, Equation~\ref{eq:pr}), which we call a \emph{concentrated} circuit, and DFC on $254$, a \emph{distributed} one. The two are matched on test accuracy and on the mean gap between the correct logit and the largest incorrect one, which we call the \emph{logit margin}, so neither circuit simply fits the data better than the other (Table~\ref{tab:twoaxis}, Figure~\ref{fig:endpoints}).

The computation itself is the same. Both rules find the periodic Fourier algorithm that existing theory derives from the task and the architecture \citep{nanda2023progress,gromov2023grokking}: keeping only the strongest input frequencies recovers $95\%$ of accuracy for either, neither builds the alternative algorithm reported by \citet{zhong2023clock}, their pizza construction, and changing the rule costs no more functional agreement than changing the seed (Appendix~\ref{app:algo}). That DFC groks at all is, to our knowledge, the first demonstration of grokking under a control driven rule.

\subsection{The credit knob}
\label{sec:rules}
Backpropagation and DFC differ in many things at once, so the gap between $154$ units and $254$ does not yet show that the credit split caused it. We therefore set DFC aside and work inside backpropagation alone, changing how evenly its credit is spread across units and reading off the resulting circuit. Everything else is held fixed: the same network, data, initialization, optimizer, learning rate, loss, and weight decay of $10^{-3}$, a setting backpropagation needs to grok at all \citep{power2022grokking,nanda2023progress}. The optimizer and the learning rate are shared, so the weight decay acts the same way in every run, and a structural difference can come only from how credit reaches the weights. The instrument is the \emph{credit knob}, built to ablate the credit split rather than to imitate a published rule.

The knob interpolates backpropagation's per-unit credit toward uniform with strength $\beta \in [0,1]$. The gradient of a weight matrix has one row per unit of the layer it feeds, row $j$ holding the gradient of every weight going into unit $j$. Writing $\vg_j$ for that row, so that $\|\vg_j\|$ is the credit unit $j$ receives, and $\bar{g} = \frac{1}{n}\sum_l \|\vg_l\|$ for the mean row norm,
\begin{equation}
\vg_j \;\leftarrow\; \frac{\vg_j}{\|\vg_j\|} \Big( (1-\beta)\, \|\vg_j\| + \beta\,
\bar{g} \Big),
\label{eq:knob}
\end{equation}
which preserves every row's direction and changes only how magnitude is spread across units. It rescales the weights entering a hidden layer, one unit at a time, and those units are the ones circuit participation is read from, so credit and circuit describe one set of neurons from opposite sides. The effect stays where it is applied, since equalizing one layer redistributes that layer and leaves the others untouched (Appendix~\ref{app:imposed}), and the transformer runs apply the knob to every linear map (Appendix~\ref{app:transformer}). At $\beta = 0$ the knob is plain backpropagation, and at $\beta = 1$ every unit receives credit of equal magnitude. The rescaled update is a descent direction at every $\beta$ (Proposition~\ref{prop:descent}) and still minimizes backpropagation's cross entropy loss, so the two ends are loss matched by construction rather than by tuning, and Appendix~\ref{app:baselines} measures what alternative interventions change instead.
Since the knob changes no term of the objective, the set of minima is the same at every $\beta$ and only which of them the run ends in changes.

Equation~\ref{eq:pr} on the row norms $\|\vg_j\|$ says how evenly a rule spreads credit, the quantity a rule sets directly, and we call it \emph{credit participation}, credit PR beside circuit PR in tables. The two describe those neurons from either side: circuit participation reads their outgoing weights in the trained network, credit participation the incoming updates a rule delivers while it trains, taken after that rule's own rescaling and averaged over the last $800$ epochs. The vector of per-unit credit magnitudes $(\|\vg_j\|)_j$ is a rule's \emph{credit profile}.

\paragraph{The credit split decides how many units the circuit uses.} Turning the knob from $\beta=0$ to $\beta=1$ equalizes credit and drives circuit participation from $0.60$ to $0.99$ over ten seeds, past where DFC sits, at fixed loss, margin, norm and test accuracy (Figure~\ref{fig:steer}). Since nothing else varied, how evenly credit is split across units is what decides how many units carry the computation, and Proposition~\ref{prop:select} works the mechanism through in a model small enough to solve exactly. Two Fourier readouts rise with it: \emph{spectral entropy}, how widely the input weights' Fourier power \citep{nanda2023progress} is spread over frequencies, and \emph{per neuron periodicity}, the mean share of a unit's power on its strongest single frequency. Both are invariant to output scale and written out in Appendix~\ref{app:repro}. The knob is not simply breaking the network: both ends reach test accuracy above $0.97$, margins between $4.94$ and $5.16$ and training cross entropy $0.040$, so it selects among grokked solutions. Full equalization costs about two points of test accuracy, which $\beta=0.25$ avoids while still reaching a highly distributed circuit.

Credit and circuit are tightly coupled here: a unit's credit and its outgoing norm correlate at $0.995$ across units. A correlation that size does not say which of the two drives the other, and both could instead be tracking how much each unit matters to the task. We therefore remove, one at a time, three things the credit could be working through: the information in the profile, its persistence, meaning whether the same units keep receiving the same credit, and its timing, meaning when during training it acts (Appendix~\ref{app:imposed}). \emph{An imposed profile strips the information.} An artificial profile, drawn to a chosen participation and dealt to units at random, makes the circuit match its shape to a mean absolute deviation of $0.018$ across five targets, though it carries no information about any unit's role. \emph{A permuted profile strips the persistence.} Shuffling backpropagation's own magnitudes across units at every step destroys the concentration, landing at $0.992$, so a circuit concentrates only if credit focuses on the \emph{same} units over time. \emph{A switch after grokking strips the timing.} Turning the knob on only after backpropagation has grokked reorganizes the finished circuit from $0.61$ to $0.99$ within five thousand epochs at test accuracy $0.997$, and turning it off re-concentrates it, so the layout is a state the rule maintains rather than one locked in at the transition. Given the three interventions, the per-unit credit profile is sufficient to set the circuit's concentration. The gap between backpropagation and DFC is therefore a difference in credit, not two separate phenomena. Both ends hold their structure for $200{,}000$ epochs, ten times our standard budget, so they are two distinct solutions rather than one reached at two speeds (Appendix~\ref{app:long}).

\subsection{Grain is the mechanism, and it places published optimizers}
\label{sec:law}

The knob shows what per-unit credit does when we set it ourselves. Published optimizers rescale credit as well, each over its own group of parameters, and that group is what decides whether the circuit moves at all. The eight optimizers of Section~\ref{sec:ruleset} differ in which parameters share one normalizer, and we call that group of parameters the rule's \emph{grain}. The knob's grain is one unit, and to vary the grain directly we generalize Equation~\ref{eq:knob} at $\beta = 1$ to a partition of the $n$ rows into groups of $m$ units, rescaling all rows of a group by one common factor so that every group ends up with the same mean credit:
\begin{equation}
\vg_j \;\leftarrow\; \vg_j\, \frac{\bar{g}}{\bar{g}_{B(j)}},
\qquad
\bar{g}_{B(j)} \;=\; \frac{1}{m}\sum_{l \in B(j)} \|\vg_l\|,
\label{eq:grain}
\end{equation}
where $B(j)$ is the group containing unit $j$ and $\bar{g}$ is the overall mean row norm of Equation~\ref{eq:knob}. At $m = 1$ this is the knob at $\beta = 1$, and at $m = n$ it rescales the whole layer at once and changes nothing.

A single factor multiplies every unit in its group by the same amount, so it can shift credit from one group to another but cannot change how credit is divided \emph{within} a group. If one unit in a group starts with twice the credit of another, it keeps twice the credit after any common factor is applied. Uneven credit inside a group therefore survives any amount of equalizing, and it disappears only once each group holds a single unit. Proposition~\ref{prop:grain} derives the prediction that follows: circuit participation stays near backpropagation's value while groups hold several units, and reaches the distributed end only once every group holds one.

We test this by running Equation~\ref{eq:grain} at shrinking group sizes. From whole-layer groups down to eight units per group the circuit does not move: it stays within $0.023$ of the circuit participation backpropagation reaches when no credit is equalized at all, which we call \emph{plain flow}. It reorganizes only over the last three halvings, from eight units per group to four, two and one, and the same threshold appears at $p=31$, $p=47$, $p=59$ and on $S_5$ (Figure~\ref{fig:law}b plots circuit participation against group size). Proposition~\ref{prop:grain} predicts the whole curve to within $0.08$ from a single measurement at the whole-layer setting, with nothing else fitted (Appendix~\ref{app:grain}).

\begin{figure}[tb]
\vspace{-1mm}\begin{center}\includegraphics[width=0.82\linewidth]{figures/fig_law.pdf}\end{center}\vspace{-3mm}
\caption{Grain decides which of the two branches a rule falls into. Every marker is measured, every rule is labelled, and the two lines are the prediction, fitted to nothing: a per-unit rule, or one that rescales nothing, should land on circuit $=$ credit, its circuit as spread as its credit, and any other grain at plain flow, backpropagation's own $0.60$. \textbf{(a)} Circuit against credit participation for the grain sweep, equalizing credit across groups of units at $25$ settings on four tasks, with the named rules of Table~\ref{tab:twoaxis} (stars) and the six optimizers that normalize per parameter, per matrix or factored (triangles). \textbf{(b)} The same sweep against group size, where the circuit moves only once groups reach single units. Bars one standard deviation, three seeds.}
\label{fig:law}
\end{figure}

\paragraph{The grain law.} \emph{Only a rule that rescales credit one unit at a time can move the circuit at all. When it does, the circuit participation comes out equal to the credit participation, so even credit gives a distributed circuit and uneven credit a concentrated one. Any other grain leaves the circuit concentrated, whatever its credit looks like.} The law follows from Proposition~\ref{prop:grain}'s model and can be applied to a rule by inspecting its update equation, without training it. Backpropagation and the knob are the two cases already seen, one rescaling nothing and staying concentrated, the other rescaling per unit and following the credit it applies. DFC belongs with the knob, though we had to measure it rather than read it off an equation, since its credit is not a transformation of the gradient yet still reaches each unit separately. Across grain settings the mean gap between credit participation and circuit participation is $0.0074$ within a seed, the vertical distance from the circuit $=$ credit line in Figure~\ref{fig:law}a, while at every other grain the circuit ignores the credit and stays where plain flow would have left it.

The published optimizers are the second comparison: credit stays untouched and the whole optimizer is swapped, each rule at its own tuned learning rate. Weight decay is then not strictly comparable across rules, so we use these runs only to read off how many units each rule's circuit uses. RMSProp, LAMB, Adafactor, Lion, Muon and Shampoo all have grains other than one unit, so all six are predicted concentrated, and all six grok and land between $0.470$ and $0.555$ (Table~\ref{tab:zoo}, full scoring in Appendix~\ref{app:zoo}). LARS normalizes exactly as the knob does but per layer, and it stays concentrated too, at $0.613$. Nero is the set's one per-unit rule, so it alone is predicted distributed, and it lands there at $0.986$ (Appendix~\ref{app:named}). Lion is the sharpest test of the two. Its grain is per parameter, which predicts a concentrated circuit, and that is where it lands, at $0.535$. Its credit, on the other hand, is perfectly even, at credit participation $1.000$, because a sign update gives every row the same norm, so reading the credit alone would have placed it at the distributed end. Grain is the more accurate rule in general, not only here: predicting each rule from its credit fails four times in eight, where grain is right eight times in eight.
